% Auto-generated from meeting_2567-01-19_clo-development.md (กบ.วช. format v2)
\documentclass{meetingminutes}
\meetingcommittee{คณะกรรมการบริหารหลักสูตร วิศวกรรมคอมพิวเตอร์และปัญญาประดิษฐ์ (บัณฑิตศึกษา)}
\meetingnumber{๑}{๒๕๖๗}
\meetingdate{วันที่ ๑๙ มกราคม พ.ศ. ๒๕๖๗}
\meetingplace{ห้องประชุมคณะเทคโนโลยีอุตสาหกรรม}
\meetingstart{๐๙.๐๐ น.}
\meetingend{๑๑.๑๕ น.}

\begin{document}
\maketitle

\psection{ผู้มาประชุม}
\begin{participants}
  \pmember{1}{รองศาสตราจารย์ ดร.อิสระ อินจันทร์}{ที่ปรึกษาหลักสูตร}{ประธาน}
  \pmember{2}{รองศาสตราจารย์ ดร.กันต์ อินทุวงศ์}{ที่ปรึกษาหลักสูตร}{รองประธาน}
  \pmember{3}{ผู้ช่วยศาสตราจารย์ ดร.พิศิษฐ์ นาคใจ}{อาจารย์ประจำหลักสูตร}{กรรมการ}
  \pmember{4}{ผู้ช่วยศาสตราจารย์ ดร.กาญจนา ดาวเด่น}{อาจารย์ประจำหลักสูตร (เลขานุการ)}{กรรมการ/เลขานุการ}
  \pmember{5}{ผู้ช่วยศาสตราจารย์ ดร.พัชรี มณีรัตน์}{อาจารย์ประจำหลักสูตร}{กรรมการ}
  \pmember{6}{ผู้ช่วยศาสตราจารย์ ดร.พิทักษ์ คล้ายชม}{รองคณบดีฝ่ายวิชาการ คณะเทคโนโลยีอุตสาหกรรม}{กรรมการ}
  \pmember{7}{ผู้ช่วยศาสตราจารย์ ดร.โสกณ วิริยะรัตนกุล}{อาจารย์ประจำหลักสูตร}{กรรมการ}
  \pmember{8}{ผู้ช่วยศาสตราจารย์ ดร.ธัชชัย อยู่ยิ่ง}{อาจารย์ประจำหลักสูตร}{กรรมการ}
  \pmember{9}{ผู้ช่วยศาสตราจารย์ ดร.อภิศักดิ์ พรหมฝ่าย}{อาจารย์ประจำหลักสูตร}{กรรมการ}
\end{participants}

\psection{ผู้ไม่มาประชุม}
\noindent ไม่มี\par

\startmeeting
\openingchair{รองศาสตราจารย์ ดร.อิสระ อินจันทร์}{กล่าวเปิดการประชุมและดำเนินการประชุมตามระเบียบวาระ ดังนี้}

\agenda{๑}{รับรองรายงานการประชุมครั้งที่ 4/2566}
ที่ประชุมรับรองรายงานการประชุมครั้งที่ 4/2566 โดยไม่มีการแก้ไข


\agenda{๒}{พิจารณา CLO ของวิชาเฉพาะด้านบังคับ}
อาจารย์แต่ละท่านนำเสนอ CLO ของวิชาที่รับผิดชอบ:

วิชา 4095101 — ปัญญาประดิษฐ์และการเรียนรู้ของเครื่อง ผู้สอน: ผศ.ดร.พิศิษฐ์ นาคใจ + ผศ.ดร.พัชรี มณีรัตน์ - CLO1: อธิบายหลักการและทฤษฎีของปัญญาประดิษฐ์และการเรียนรู้ของเครื่องได้ (Bloom Cog 2) → PLO1 - CLO2: ใช้เครื่องมือ AI/ML ออกแบบและพัฒนาโมเดลตามบริบทที่กำหนดได้ (Bloom Cog 3) → PLO1 - CLO3: ประเมินประสิทธิภาพโมเดล AI/ML และเลือกแนวทางที่เหมาะสมได้ (Bloom Cog 4) → PLO3 - CLO4: ทำงานเป็นทีมและสื่อสารผลการวิเคราะห์ได้อย่างถูกต้อง (Bloom Aff 3) → PLO5 - CLO5: ปฏิบัติตามจรรยาบรรณในการใช้ AI อย่างมีความรับผิดชอบ (Bloom Psy 3) → PLO4

วิชา 4095102 — การเขียนโปรแกรมขั้นสูงสำหรับการเรียนรู้ของเครื่อง ผู้สอน: ผศ.ดร.พิศิษฐ์ นาคใจ + ผศ.ดร.พัชรี มณีรัตน์ - CLO1: เขียนโปรแกรมขั้นสูงด้วย Python/PyTorch/TensorFlow สำหรับ ML ได้ → PLO1 - CLO2: ออกแบบ ML Pipeline และ Workflow ที่มีประสิทธิภาพได้ → PLO1, PLO2 - CLO3: บูรณาการ Git/GitHub และ MLOps ในโครงงานจริงได้ → PLO1, PLO3 - CLO4: สร้างผลงาน ML ที่นำไปใช้แก้ปัญหาจริงในองค์กรหรือชุมชนได้ → PLO2, PLO4

ที่ประชุมอภิปรายและเห็นชอบ CLO ของวิชาทั้ง 2 ดังกล่าว โดยยืนยันว่าระดับ Bloom's สอดคล้องกับระดับ PLO


\agenda{๓}{จัดทำตาราง Curriculum Mapping (PLO x CLO)}
ที่ประชุมร่วมกันสร้างตาราง Curriculum Mapping แสดงความสัมพันธ์ระหว่าง CLO ของทุกวิชากับ PLO ทั้ง 5 ข้อ โดยกำหนดระดับความสัมพันธ์เป็น Primary (●) และ Secondary (○)

\resolution{เห็นชอบตาราง Curriculum Mapping เบื้องต้น และมอบหมายให้อาจารย์แต่ละท่านพัฒนา CLO ของวิชาเลือกที่รับผิดชอบ เพื่อนำเสนอในการประชุมครั้งถัดไป}

\agenda{๔}{ความก้าวหน้าการจัดทำ มคอ.2}
ผศ.ดร.พิศิษฐ์ รายงานว่า มคอ.2 อยู่ระหว่างการจัดทำ คาดว่าจะเสร็จสมบูรณ์ภายในเดือนกุมภาพันธ์ 2567 เพื่อเสนอต่อคณะกรรมการบัณฑิตศึกษาต่อไป


\adjournmeeting

\begin{sigblock}
\typist{ผู้ช่วยศาสตราจารย์ ดร.กาญจนา ดาวเด่น}
\reviewer{รองศาสตราจารย์ ดร.อิสระ อินจันทร์}{กรรมการ}
\chairsig{รองศาสตราจารย์ ดร.อิสระ อินจันทร์}{ประธานการประชุม}
\end{sigblock}

\end{document}
